\documentclass[10pt,a4paper]{amsart}
\usepackage[margin=19mm]{geometry}
\usepackage{amsmath,amssymb,mathtools}
\usepackage[scr=rsfs,cal=euler]{mathalfa}
\usepackage{xfrac}
\usepackage[unicode,hidelinks,bookmarks=false]{hyperref}
\newcommand{\ie}{i.e.\ }
\newcommand{\cf}{cf.\ }
\newcommand{\resp}{respectively\ }
\newcommand{\Subsection}{\subsection}
\providecommand{\nobreakdash}{\nobreak\hbox{-}}
\setlength{\emergencystretch}{2em}
\multlinegap0pt
\usepackage{bm}
\usepackage{bbm}
\usepackage{dsfont}
\usepackage{xspace}
\usepackage{cancel}
\usepackage{mathtools}
\usepackage[shortlabels]{enumitem}
\usepackage{amsthm}
\makeatletter
\@ifundefined{theorem}{\newtheorem{theorem}{Theorem}}{}
\@ifundefined{lemma}{\newtheorem{lemma}[theorem]{Lemma}}{}
\@ifundefined{proposition}{\newtheorem{proposition}[theorem]{Proposition}}{}
\makeatother
\DeclareRobustCommand*{\ora}{\overrightarrow}
\newcommand\restr[2]{{  \left.\kern-\nulldelimiterspace #1 \vphantom{\big|} \right|_{#2}  }}
\newcommand{\R}{\mathbb{R}}
\newcommand{\barcode}[1]{ {\mathcal B}\left(#1\right) }
\newcommand{\conv}[1]{ {\rm conv}(#1) }
\newcommand{\field}{k}
\newcommand{\induced}[1]{\field^{#1}}
\newcommand{\interior}[1]{{\kern0pt#1}^{\mathrm{o}}}
\newcommand{\lowerb}[1]{ L[#1] }
\newcommand{\norm}[1]{\left\lVert #1 \right\rVert}
\newcommand{\rebuttal}[1]{{ #1}}
\newcommand{\upperb}[1]{ U[#1] }
\title[Multi-parameter Module Approximation: an efficient and inter]{Multi-parameter Module Approximation: an efficient and interpretable invariant for multi-parameter persistence modules with guarantees}
\author{David Loiseaux, Mathieu Carri\`ere, Andrew J. Blumberg}
\date{Source: arXiv:2206.02026v4 (CC BY 4.0)}
\begin{document}
\maketitle

\noindent\emph{Source and licence.} Multi-parameter Module Approximation: an efficient and interpretable invariant for multi-parameter persistence modules with guarantees. Version arXiv:2206.02026v4 (CC BY 4.0), distributed under the Creative Commons Attribution 4.0 International licence, as recorded in the arXiv licence metadata for this exact version and shown on the abstract page https://arxiv.org/abs/2206.02026v4. Full rights notice: licenses/arxiv-2206.02026-CC-BY-4.0.txt.\par\smallskip

\noindent\textit{Editorial scope.} Counted unit: the Proposition whose source label is \texttt{prop:local\_bound} in the article (source0.tex lines 2065--2078, source label \texttt{prop:local\_bound}), delivered together with the article's own complete original proof of it (source0.tex lines 2080--2122, one \texttt{proof} environment of 43 lines). No mathematics is written, repaired or re-derived: statement and proof are the article's own text line for line. Required context retained uncounted: lemma:compatible\_close (lines 1011--1031); lemma:boundary\_bound (lines 1977--1991); eq:surr (lines 2039--2041). Every definition, display and label of those lines is retained unchanged. Omitted from this package and not redistributed: 1--1010, 1032--1976, 1992--2038, 2042--2064, 2079--2079, 2123--3745. Nothing inside a retained excerpt is omitted or altered: every retained body is delivered exactly as the article's lines are written, so no omission receipt of any kind accompanies this package. Citations: the retained excerpts carry no citation command at all, so no bibliography entry of the article is transcribed and no reference is redistributed. Reference adaptation: none was needed and none was made; every reference inside the delivered unit resolves inside it, so no unresolved reference or citation is produced. Printed slips kept literal: no printed word, symbol, hypothesis or number of the counted statement or of its proof was altered. Layout and class: the article's own class is \texttt{article} with packages including libertinus, authblk, url, etoolbox, inputenc, multicol, amsmath, amssymb, amsfonts, amsthm, mathrsfs, appendix, xcolor, textcomp; the wrapper is the lane's pinned \texttt{amsart} editorial wrapper at 10pt on A4, which restates the theorem-like environments the retained text needs and adds the article's own macro definitions that the retained text needs (\texttt{\textbackslash R}, \texttt{\textbackslash barcode}, \texttt{\textbackslash conv}, \texttt{\textbackslash field}, \texttt{\textbackslash induced}, \texttt{\textbackslash interior}, \texttt{\textbackslash lowerb}, \texttt{\textbackslash norm}, \texttt{\textbackslash ora}, \texttt{\textbackslash rebuttal}, \texttt{\textbackslash restr}, \texttt{\textbackslash upperb}), with extra wrapper packages (bm, bbm, dsfont, xspace, cancel, mathtools, enumitem). The delivered numbering is the unit's own. Assets: no retained span contains a figure environment, a graphics-inclusion command, a PDF-page inclusion, a table or a TikZ picture; no image file is copied into this package, the package declares no asset dependency and no image file is redistributed with it. Licence: version arXiv:2206.02026v4 of this article is distributed under the Creative Commons Attribution 4.0 International licence, as recorded in the arXiv licence metadata for this exact version and shown on the abstract page https://arxiv.org/abs/2206.02026v4. A full rights notice is delivered at licenses/arxiv-2206.02026-CC-BY-4.0.txt. Characters: every retained excerpt is pure ASCII, so the delivered engine, XeLaTeX with its default Latin Modern text font, reports no missing character.
\begin{lemma}\label{lemma:compatible_close}
	Let $\rebuttal{k^I}$ be a \rebuttal{f.p.} interval module, let $l_1,l_2\subseteq\R^n$ be two diagonal lines and
	let $\ora u\in \R^n$ be a positive or negative vector
	such that $l_2 = l_1 + \ora u$.
	Then, the following properties hold:
	\begin{enumerate}[label=\textnormal{(\textit{\roman*})}]
		\item If the barcode
			$\barcode{\restr{\rebuttal{\field^I}}{l_1}}=\{[b^I_{l_1}, d^I_{l_1})\}$
			is not empty and satisfies $\norm{d^I_{l_1} - b^I_{l_1}}_\infty >
			\norm{\ora u}_\infty$, then the barcode
			$\barcode{\restr{\rebuttal{\field^I}}{l_2}}$ is not empty as well, and
		\item	If the barcodes $\barcode{\restr{\rebuttal{\field^I}}{l_1}}$ and
			$\barcode{\restr{\rebuttal{\field^I}}{l_2}}$ are not empty,
			then one has
			\begin{equation*}
				\norm{ d^I_{l_1}-d^I_{l_2}}_{\infty} \le \norm{ \ora u}_\infty \text{ and } \norm{ b^I_{l_1}-b^I_{l_2}}_\infty \le \norm{ \ora u}_\infty,
			\end{equation*}
			where we used the conventions
			$(+\boldsymbol\infty) - (+\boldsymbol\infty) = (-\boldsymbol\infty) - (-\boldsymbol\infty)  = 0 $.
	\end{enumerate}
\end{lemma}

\begin{lemma}[Endpoints bound]\label{lemma:boundary_bound}
	Let $k^I$ be a {f.p.} interval module. Let $K$ be a rectangle in $\R^n$ and $L:=L_\delta(K^\delta)$ be the $\delta$-grid of lines of the offset $K^\delta$.
	Let $x\in K$, $l_x$ be the diagonal line passing through $x$, and let $
	L_{x,\delta}:=\{l\in L\,:\, d_\infty(x,l) \leq \delta \text{ and } l_x, l \text{ are $\delta$-comparable}\}
	$, which is non-empty since
	$L$ $\delta$-fills $K$.
	Assume that $l_x\cap I\neq \varnothing$, and $d^I_x \in \upperb{I}$ be the associated deathpoint, and
	assume that for any line $l$ in $L_{x,\delta}$, one has $I\cap l\neq \varnothing$, and
	let $D_{x,\delta}^I$ be the set of the associated deathpoints: $D_{x,\delta}^I=\{d_l^I\,:\, l\in L_{x,\delta}\}$.
	Then, $d^I_x$ belongs to the rectangle hull of a subset $\tilde
	D_{x,\delta}^I$ of $D_{x,\delta}^I$: one has $d^I_x\in \mathrm{recthull} [\tilde D_{x,\delta}^I]$ with $\tilde D_{x,\delta}^I \subseteq D_{x,\delta}^I$.

	Similarly, if $b^I_x\in \lowerb{I}$ is a birthpoint, then $b^I_x\in\mathrm{
	recthull}[\tilde B_{x,\delta}^I]$, where $\tilde B_{x,\delta}^I$ is a subset of $B_{x,\delta}^I=\{b_l^I\,:\, l\in L_{x,\delta}\}$, i.e., the set of birthpoints associated to $L_{x,\delta}$. \\
\end{lemma}

	\begin{equation}\label{eq:surr}
		\Vert d^+ - d\Vert_\infty, \Vert d^- - d\Vert_\infty \le \delta.
	\end{equation}

\begin{proposition}\label{prop:local_bound}
	Let ${k^{I}}$ be a {f.p.} interval module. Let $K$ be a rectangle in $\R^n$ and $L:=L_\delta(K^\delta)$ be the $\delta$-grid of lines of the offset $K^\delta$.
	Let $l\in L$ such that $|L_l| = 2^{n-1}$ and assume that $\conv{L_l} \cap \lowerb{I}$ (resp. $\conv{L_l} \cap \upperb{I}$) is not empty.
	Then, there exists a set $B_l$ (resp. $D_l$) such that for any $x\in \conv{L_l} \cap \lowerb{I}$ (resp. $\conv{L_l} \cap \upperb{I}$),
	one has either  $\Vert b_x^I -d_x^I \Vert_\infty \le \delta$, or
	$x \in B_l$ (resp. $D_l$), where $B_l$ (resp. $D_l$) is a rectangular set in
	$\R^n$ that can be constructed from the birthpoints $(b_{l'}^I)_{l'\in L_l}$  (resp. deathpoints $(d_{l'}^I)_{l'\in L_l}$). %
	Moreover, one has that:
	\begin{equation}\label{eq:diam}
		{\rm diam}(B_l)={\rm sup}_{x,x'\in B_l}\norm{x-x'}_\infty \leq \delta,
	\end{equation}
	and similarly for $D_l$. %

\end{proposition}

\begin{proof} We first construct $B_l$ and $D_l$, and then we will show items (1) and (2). \\

	{\bf Definition of $B_l,D_l$.}
	Let first assume that $x$ is in the interior of $ \conv{L_l}$, that we denote with $\interior{\conv{L_l}}$.
	Note that if there is a line $l_0$ that is $\delta$-comparable to $l_x$, and such that $\barcode{\restr{\induced{I}}{l_0}} = \varnothing$,  %
	then by Lemma~\ref{lemma:compatible_close} $(i)$, one immediately has $\Vert b_x^I - d_x^I \Vert_{\infty} \le \delta$.
	Hence, we now assume that the barcodes along any %
	line that is $\delta$-comparable to $l_x$ is not empty,
	which means that the hypotheses
	of Lemma~\ref{lemma:boundary_bound} are satisfied for $x$.
	Now, remark that since $L$ is a grid, if one is able to find a line $l'$ in $L$ whose intersections
	with hyperplanes associated to the canonical axes of $\R^n$ are $\delta$-close to $x$, then, since $x$
	is in the interior of an $L$-surrounding set $L_l$, $l'$ must belong to that surrounding set $L_l$ as well. More formally,
	one has that, for any line $l'\in L$:
	\begin{equation*}
		d_{\infty}(x, l'\cap H_i)\le \delta \quad\Longrightarrow\quad l' \in L_l,
		\quad
		\text{ where }
		H_i = \left\{ y \in \R^n : y_i = x_i\right\}.
	\end{equation*}

	This ensures (from Equation~\ref{eq:surr}) that the lines of $L$ associated to $\tilde B^I_{x,\delta}$ and  $\tilde D^I_{x,\delta}$
	are all included in $L_l$ for any $x\in\interior{\conv{L_l}}$, and thus that we can safely define:
	\begin{equation*}
		D_l := \overline{\bigcup_{x \in \interior{\conv{L_l}}}\mathrm{ recthull}[\tilde D_{x,\delta}^I]}
		\quad \text{and} \quad
		B_l := \overline{\bigcup_{x \in \interior{\conv{L_l}}}\mathrm{ recthull}[\tilde B_{x,\delta}^I] }.
	\end{equation*}
	Note that $B_l$ and $D_l$ depend only on the endpoints of the lines in $L_l$ %
	and that $d_x^I \in D_l$ and $b_x^I \in B_l$ for any $x\in \interior{\conv{L_l}}$ by Lemma~\ref{lemma:boundary_bound}.
	Furthermore, if $x$ is in the closure of $\conv{L_l}$, the previous statements still hold
	since $D_l$ and $B_l$ are closed sets.
	We now show that $B_l$ and $D_l$ satisfy Equation~(\ref{eq:diam}). \\

	{\bf Proof of Equation~(\ref{eq:diam}). } By applying
	Lemma~\ref{lemma:boundary_bound} and its proof for arbitrary dimension $j$ to all $x\in\conv{L_l}\cap U[I]$, %
	there exist deathpoints $\overline d^j$ and $\underline d_j$
	that satisfy $(\overline d^j)_j = \sup_{d\in D_l} d_j$ and $(\underline d_j)_j = \inf_{d\in D_l} d_j$ and $(\underline d_j)_j \le x_j \le (\overline d^j)_j$ for all $x\in\conv{L_l}\cap U[I]$.
	Moreover, these points are located on lines in $L_l$
	that are $\delta$-consecutive by definition.
	Thus, applying Lemma~\ref{lemma:compatible_close} $(ii)$ repeatedly on these
	pairs of lines for all dimensions, we end up with $D_l$ having a diagonal smaller than $\delta$. The same goes for birthpoints. %
\end{proof}

\end{document}
